\chapter{环境对象的结构化表示：测量、关系与约束}
\label{chp:physics_anatomy}

我们建立智能理论，是为了说明有限系统怎样组织信息、更新状态并形成有效行动。环境是这些过程的约束来源，却不是内部表示的现成副本。一个苹果可以被拍摄、称量、分类和抓取；不同任务需要不同的对象属性、空间关系与预测变量。我们不必在系统中重建苹果的全部微观构成，才能完成这些任务；也不能因为系统生成了苹果的图像，就认定它掌握了苹果在所有条件下的行为。本章从测量开始，依次建立环境状态、观测记录、结构化估计和任务预测之间的关系。涉及位置、力或电量时，我们保留物理单位和领域规律，因为这些知识约束所描述的对象。内部编码、信息传播与判断更新则使用计算动力学语言。我们由此区分对象的变化与模型的变化，并给出检验二者对应程度的方法，使环境知识能够进入一般智能理论，而不成为关于认知机制的未经证明的普遍断言。

\section{从环境状态到观测记录：测量不是复制}
\label{sec:general_relativity}

我们首先区分环境状态与观测记录。环境状态包含足以描述目标过程的变量；观测记录则是传感器、采样程序或数据接口实际返回的结果。对苹果进行称量时，秤的读数并不是苹果全部状态，而是某种测量条件下的质量估计。拍照时，像素受到视角、照明、遮挡与设备响应的共同影响。系统把像素变成对象属性，需要经过对象定位、特征提取与解释选择。这些环节都会引入误差，也会舍弃与当前任务无关的信息。把观测直接命名为真实状态，会掩盖这些中间环节，使后续推理获得不应有的确定性。

令环境状态为 $z_k\in\mathcal Z$，观测为 $o_k\in\mathcal O$，动作或测量设置为 $a_k$。我们采用下述一般接口：
\begin{equation}
z_{k+1}=T(z_k,a_k,\eta_k),\qquad o_k=H(z_k,a_k,\xi_k).
\end{equation}
其中两个映射可以是确定性的，也可以通过随机变量描述未建模因素。这里的随机性不必表示对象本身具有某种特殊性质；它也可以表达我们没有记录的条件。若需要概率推断，就必须说明噪声分布与条件依赖。若只有有限样本而没有可信的概率模型，我们可以报告经验误差和适用区间，却不能把任意评分直接解释成精确概率。状态空间的选择同样属于建模工作，而非对环境本体的完整宣告。

观测记录应携带来源、时间、单位和测量设置。两个温度数值只有经过单位转换、时间对齐和对象匹配，才可以进入同一次比较。来自不同传感器的坐标还需要坐标变换；变换本身的误差也应保存。无量纲特征向量可以用于检索，但不能与长度、质量等测量值未经转换便逐项相加。我们因此把内容分成具有明确量纲的测量字段、经过规定变换的编码字段，以及关于可信程度的元数据。三者可以共同参与决策，却承担不同的解释职责。

对象身份也是测量的一部分。连续两帧中出现两个相似苹果，并不意味着系统已经知道哪个苹果对应上一帧的目标。身份匹配可以利用位置、外观、运动预测和任务上下文，但这些依据可能冲突。我们保存候选对应及其支持证据，而不是提前把所有候选合并为一个对象。若随后发现目标被替换，应修订身份关系及其派生预测，不能只更新颜色字段而让错误身份继续影响行动。结构化记录的价值正在于把这种关联错误暴露出来。

缺失观测必须与零值区分。一个没有称量过的苹果，其质量未知；它并不因此具有零质量。我们给属性添加有效性标记，规定哪些操作允许使用估计、哪些操作必须等待测量。遮挡期间可以依据模型维持预测，但预测记录需要显示其起点和经过的时间。随着缺少校正的时间增长，系统可能需要扩大误差区间或重新定位。持续维持一个数值，并不能证明对对象的追踪仍然正确。

观测设计还决定哪些问题能够得到回答。如果传感器只记录外观，系统可能无法区分外表相似但内部状态不同的两个苹果；增加称量、光谱或接触测量后，原有的等价候选才可能被区分。我们把这种不可区分性理解为接口与任务的限制，不把它简单归结为编码维度不足。增加向量长度而不增加有效信息，并不会自动解决测量不足。对 AgentOS 而言，这意味着工具调用应以待消除的不确定性为依据：先说明当前决策需要辨别什么，再选择能够提供相应证据的观测动作。

\section{从观测到联合估计：关系与属性共同消歧}
\label{sec:quantum_field}

我们用联合估计 $(S_k,X_k)$ 表示系统当前对环境的组织方式。结构部分记录对象、区域、事件及其类型化关系；属性部分记录测量值、类别编码和置信信息。这种分工不要求两个部分分别从输入中独立提取。苹果是否位于盘中，会影响区域归属；盘沿的轮廓又可能帮助判断苹果是否被遮挡。若分别做出互不协调的判断，再机械拼接结果，模型就可能同时声称苹果被盘沿完全遮挡且其完整边界已经直接可见。我们需要让结构与属性在相同的对象身份和证据条件下接受联合检查。

构造过程可以从候选对象开始。候选不是已经确证的实体，而是输入中可能需要独立解释的一组特征。系统为候选指定临时身份，提出类别与关系假设，然后计算这些假设对观测的解释程度。若采用概率模型，令 $h$ 表示包含结构和属性的一个候选解释，给定观测 $o$，在有限候选集上可写成
\begin{equation}
p(h\mid o)=\frac{p(o\mid h)p(h)}{\sum_{h'}p(o\mid h')p(h')},
\end{equation}
前提是分母为正且似然与先验已经定义。该公式来自条件概率的归一化，并不保证候选集覆盖真实情况。所有候选都很差时，最高后验候选仍可能错误。我们因此保留拒绝解释或扩展候选集的入口，而不强迫系统永远选择现有类别。

结构约束可以排除某些解释，但约束需要符合对象类型。例如，一个已经确认的单件物体在同一时刻不能被赋予两个互斥的身份位置，除非记录表达的是候选分支；一个集合则允许包含多个成员。关系中的方向、时间和作用角色必须明确。苹果在盘中与盘在苹果中不是同一断言。把包含关系错误地当成对称邻接，会使后续抓取位置与对象归属一并出错。类型化关系使我们能够在进入数值更新前检查这种错误。

我们还区分直接观测、规则推导与模型预测。图像提供了可见区域，空间重建给出了可能深度，历史轨迹预测了遮挡后的当前位置。三个字段可能都表现为数值，却不能拥有相同的证据地位。每个派生属性应指向其输入记录和使用的模型版本。当相机标定改变或源记录失效，依赖这些条件的派生结果需要重新计算。仅仅保存最终数值，将使系统无法解释为什么旧预测需要撤回，也无法确定一次修订影响哪些行动计划。

联合估计可以通过交替更新实现，但交替本身没有自动收敛保证。固定结构后调整属性，再依据新属性修订关系，可能不断在两个候选之间切换。若我们明确给出同一个目标函数，并保证每一步不增加该目标，就能得到目标值单调不增这一有限结论；仍不能据此认定状态唯一、推断正确或全局最优。对象数目改变时，目标维度和比较方式也必须说明。我们通过版本化候选、更新停止条件与冲突记录，让这种过程可复核。

苹果例子展示了联合修订的必要性。第一次观测可能把高光区域识别为独立白色物体；后续视角表明它随苹果表面共同移动，系统便可以把它修订为照明造成的外观变化。此次更新不只是删去一个节点，还要修订节点所属的属性、与盘的关系以及此前生成的对象计数。我们评价联合表示时，应检查这种依赖传播是否完成，而不只检查最后一句描述是否流畅。对于任务系统，正确的更新链条比一个孤立的漂亮输出更能说明模型已经利用环境证据。

\begin{figure}[htbp]
\centering
\includegraphics[width=0.82\linewidth,height=0.54\textheight,keepaspectratio]{image/real_apple.png}
\caption{环境对象的观察入口：苹果图像提供外观证据，关系、内部状态与行动结果仍需分别估计。}
\end{figure}

\section{预测、尺度与粗粒化：模型只保留任务所需信息}
\label{sec:condensed_matter}

结构化估计的用途不是建立一份无限细致的对象清单，而是支持指定预测与行动。辨认苹果种类时，颜色和形状可能足够；规划抓取时，还需要位置、接触条件和可用动作；判断长期保存状态时，则需要不同的观测和时间尺度。我们因此先规定任务输出，再讨论表示需要保留哪些信息。同一个环境可以产生多个合适的模型，这些模型不必具有相同维数，也不必以相同关系为核心。

令任务目标为 $y_k$，联合表示为 $x_k=e(o_{0:k},a_{0:k-1})$，读出为 $\widehat y_k=r(x_k)$。经验风险是在指定数据集合上计算的预测损失，而总体风险还依赖任务分布。即使两者的表达形式相似，也不能未经检验便相互替代。我们用训练外的数据、不同观测条件和干预试验检查模型是否保留了所需信息。若目标随任务改变，原先有效的编码可能不再充分。任务相关性不是编码的永久身份，而是编码、读出与评估条件之间的关系。

粗粒化是把多个细节状态映射成一个较简约状态。设 $c:\mathcal Z\to\mathcal C$ 为压缩映射。要使压缩状态独立支持确定性更新，必须存在 $\overline T$ 满足
\begin{equation}
c(T(z,a))=\overline T(c(z),a).
\end{equation}
若两个环境状态被压成同一个状态，但同一动作产生不同的下一压缩状态，便无法在当前压缩变量上定义唯一的更新函数。这个结论直接来自函数的单值性，不需要引入物理类比。遇到这种情况，我们可以细化状态、增加历史记忆或使用条件分布；不能只是把原来的单值预测写得更复杂。

例如，只记录苹果当前位置，可能不足以预测下一时刻位置，因为相同位置可能对应不同速度。加入速度后，在某些模型假设下预测更充分；若动作与接触状态仍未记录，模型还会出现系统性误差。这里的速度属于环境模型的领域变量，其必要性由预测任务决定。内部处理器并不因此必须具有与苹果相同的运动规律。我们把两层描述分开，才可以在保留真实领域约束的同时，为认知更新选择合适的计算实现。

尺度变化也会改变属性的含义。单个对象的局部特征、一个对象集合的统计量与整个场景的任务状态不是同一类变量。集合的平均颜色不能替代每个成员的颜色；场景中的可通行性也不能还原成某个物体的单独属性。我们通过聚合算子明确这些变量怎样形成，并记录聚合所丢失的信息。若使用平均值，应说明权重与样本范围；若使用最大值，应说明它是否代表最危险情况还是最明显证据。聚合方法本身会影响后续决策。

预测误差需要按照用途分解。身份匹配错误、测量误差、压缩损失和转移模型偏差可能产生相同的输出偏差，却需要不同修正。重新训练读出不能必然修复错误身份；增加采样频率也不能消除不合适的状态定义。我们为评估设置分层诊断：固定真实身份检查属性估计，固定属性检查关系推断，再检查整个预测流程。这样才能说明一次改进来自何处，以及它能否迁移到新的任务尺度。

对一般智能而言，重要能力是根据任务重新选择尺度并知道何时细化。系统可以先用简约模型探索，再在误差达到关键阈值时增加变量；也可以在稳定阶段把重复细节聚合为摘要。细化与压缩都应保留接口和证据回溯能力。我们不把更大模型自动视为更深理解，而以任务保持、误差可控和资源可承受作为共同标准。这为后续的状态动力学提供了可以变化但仍有明确意义的模型边界。

\section{领域规律、约束验证与行动闭环}
\label{sec:mirror_table}

环境模型需要服从对应领域的知识，但知识进入模型的方式必须明确。一个机器人抓取苹果，需要考虑尺寸、位置、接触条件和设备限制；一个文本系统整理苹果资料，则可能主要需要来源一致性、术语定义和时间信息。前者可以调用机械模型，后者可以调用文档检查规则。我们不以同一套物理方程解释两个系统的全部内部活动，而是在一般状态接口中容纳各自需要的约束与更新机制。

对连续控制任务，可将可接受状态集合写为 $K=\{z:g_j(z)\leq0,\ j=1,\ldots,m\}$。约束函数可以表达设备工作范围、对象间距或其他领域条件。若模型预测下一状态位于集合中，只说明预测满足模型约束；实际状态还受到模型误差与未记录扰动影响。系统需要为误差保留裕量，并在执行后使用观测校正。把预测可行直接理解为实际必然可行，会把模型边界之外的不确定性藏入动作接口。

在简单标量约束下，我们可以清楚展示误差如何进入决策。设实际值为 $v$，预测值为 $\widehat v$，已验证的误差界为 $|v-\widehat v|\leq\varepsilon$，要求 $v\leq b$。若采用 $\widehat v\leq b-\varepsilon$，则由 $v\leq\widehat v+\varepsilon\leq b$ 得到约束满足。这一推理只在误差界有效的条件下成立。更复杂约束需要相应的敏感性界，不能把同一个标量裕量无条件套到所有变量上。我们在公式旁说明误差界来自何处、在哪些条件下失效。

行动闭环由估计、预测、选择、执行和再观测组成。选择动作前，系统检查当前对象身份、模型版本和所需资源；执行后，系统比较预期与实际记录。若差异显著，应先判断差异是测量异常、动作未执行、环境变化还是模型错误，再决定是否继续。一个指令被生成，不等于它已经执行；工具返回成功，也不一定等于预定环境结果已经出现。我们用不同事件记录这几个阶段，以免认知状态抢先宣告任务完成。

证据不充分时，主动测量本身可以成为动作。苹果是否已经被抓住，可能需要夹持反馈或新的图像；资料中的日期冲突，则可能需要读取原始记录。我们可以比较不同测量动作预计减少的不确定性和所需成本，但成本函数必须符合当前任务，不能把所有计算、时间与错误风险不加解释地压成一个数。若采用加权和，权重表达决策偏好，其来源和调整规则也应成为模型的一部分。

领域知识还需要版本与适用条件。机械模型可能只在某个速度范围和接触类型下有效；分类模型可能只覆盖规定类别；文档规则可能依据某个时间版本。内部系统应保存这些条件，使超出范围的输入触发重新评估。通过这种机制，约束不是一段永远正确的声明，而是与数据、动作和验证流程相连的可执行知识。发现违反条件时，系统可以收缩行动范围、追加观测或切换模型。

本章所建立的环境对应关系，最终可以进入 AgentOS 的对象记录、工具接口与事件日志。HSF-HD 在这里规定一般的状态组织和更新要求，而具体传感器、机械方程或文档工具属于实现层。我们评价一个实现时，应同时检查任务效果、证据可回溯性以及更新规则是否在所声明条件下运行。环境可以非常复杂，但有效智能并不需要先宣称已经复制整个环境；它需要形成足够可靠、能够被校正且能支持行动的局部模型。
